\documentclass[10pt, oneside]{article}

% Packages
\usepackage{amsmath, amsthm, amssymb, calrsfs, wasysym, verbatim, bbm, color, graphics, geometry}
\usepackage[utf8]{inputenc}

% Geometry
\geometry{
	tmargin=.75in,
	bmargin=.75in,
	lmargin=.75in,
	rmargin=.75in
}

% Commands
\newcommand{\R}{\mathbb{R}}
\newcommand{\C}{\mathbb{C}}
\newcommand{\Z}{\mathbb{Z}}
\newcommand{\N}{\mathbb{N}}
\newcommand{\Q}{\mathbb{Q}}

% Theorems
\newtheorem{thm}{Theorem}
\newtheorem{defn}{Definition}
\newtheorem{conv}{Convention}
\newtheorem{rem}{Remark}
\newtheorem{lem}{Lemma}
\newtheorem{cor}{Corollary}

% Metadata
\title{Pregunta acerca de una posible relación entre Planitud Diferencial y formulaciones Hamiltonianas}
\author{Andres Mauricio Morales Martinez}
\date{}

\begin{document}

\maketitle

Hola profe :D

Me surgió una duda conceptual intentando conectar varias ideas del curso con cosas de mecánica analítica, así que quería preguntarle si esto tiene sentido o si estoy mezclando conceptos que no deberían mezclarse.

Entiendo que los corchetes de Lie me dan una idea geométrica de la diferencia al conmutar direcciones inducidas por campos vectoriales. Sean $\mathbf{f}$ y $\mathbf{g}$ campos vectoriales en $\mathbf{R}^n$,

\begin{equation}
	[\mathbf{f},\mathbf{g}]
	=
	\nabla \mathbf{g}\,\mathbf{f}
	-
	\nabla \mathbf{f}\,\mathbf{g}.
\end{equation}

Por otro lado, si tengo un sistema autónomo

\begin{equation}
	\dot{\mathbf{x}}(t)
	=
	\mathbf{f}(\mathbf{x}(t)),
	\label{eq:dynamic}
\end{equation}

donde $\mathbf{x}(t)$ es la curva integral del campo $\mathbf{f}$, y considero una salida escalar

\begin{equation}
	y(t)
	=
	h(\mathbf{x}(t)),
	\qquad
	h:\R^n \to \R,
\end{equation}

entonces su evolución temporal puede escribirse como

\begin{equation}
	\dot{y}(t)
	=
	\nabla h~
	\dot{\mathbf{x}}(t)
	=
	\nabla h~
	\mathbf{f}(\mathbf{x}(t))
	=
	L_f h.
\end{equation}

Hasta ahí entiendo que la derivada de Lie representa cómo evoluciona una observable cuando se sigue el flujo inducido por el campo vectorial del sistema.

Ahora, pasando al formalismo Hamiltoniano, si el sistema admite una representación Hamiltoniana con coordenadas canónicas $(q,p)$ y Hamiltoniano $H(q,p)$, la evolución temporal de una observable escalar

\begin{equation}
	\Gamma(q,p;t)
	\in
	\R
\end{equation}

está dada por

\begin{equation}
	\frac{d}{dt}
	\Gamma(q,p;t)
	=
	\{\Gamma,H\}
	+
	\frac{\partial \Gamma}{\partial t}.
\end{equation}

Si la observable no depende explícitamente del tiempo, entonces

\begin{equation}
	\dot{\Gamma}
	=
	\{\Gamma,H\}.
\end{equation}

Por lo tanto, si mi salida puede escribirse como una observable del sistema,

\begin{equation}
	y(t)
	=
	h(q(t),p(t)),
\end{equation}

parece que podría escribir

\begin{equation}
	\dot y
	=
	\{h,H\}
	=
	L_{X_H}h,
\end{equation}

donde $X_H$ sería el campo Hamiltoniano inducido por $H$. \\ 

Ahora lo interesante, es que si elijo mis estados adecuadamente de una forma que sea admitida por la formulación Hami

La intuición que me quedó es que parece existir una relación bastante íntima entre la evolución temporal escrita mediante derivadas de Lie y aquella escrita mediante corchetes de Poisson, al menos cuando la dinámica admite una representación Hamiltoniana.\\


Me surge entonces una pregunta profe, : \\



¿Tiene sentido preguntarse si existe una parametrización diferencialmente plana que además sea compatible con coordenadas canónicas? Creería inicialmente que sí y en ese caso, quizá, uno podría intentar formular el problema de encontrar una salida plana a partir de un Hamiltoniano del sistema, aprovechando la estructura geométrica asociada.

No sé si esta idea ya existe en alguna forma, si estoy mezclando conceptos incompatibles, o si realmente podría haber algo útil ahí para el método del curso, pero me dio curiosidad porque siento que la conexión entre corchetes de Lie y corchetes de Poisson podría ser más profunda de lo que estoy entendiendo ahorita.\\

Me fascinaría profundizar y explorar más en esto! No sé si me puedas ayudar. 

\end{document}